\chapter{COMPACT OPERATORS}\label{chap_cpt_ops}

\section{Definition and Elementary Properties}

We will pick up where we left off at the end of chapter 5, studying classes of operators.
Our context will be slightly more general; in the case of compact operators we will be
looking at operators on Banach spaces.  First, let us recall a few definitions and
elementary facts from advanced calculus (or a first course in analysis).

\begin{defn} A subset $A$ of a metric space $M$ is
 \index{total!boundedness}%
 \index{bounded!totally}%
\df{totally bounded} if for every $\epsilon > 0$ there exists a finite subset $F$ of $A$
such that for every $a \in A$ there is a point $x \in F$ such that $d(x,a) < \epsilon$.
For normed linear spaces this definition may be rephrased: A subset $A$ of a normed linear
space $V$ is \df{totally bounded} if for every open ball $B$ about zero in $V$ there exists
a finite subset $F$ of $A$ such that $A \subseteq F + B$.  This has a more or less standard
paraphrase: A space is totally bounded if it can be kept under surveillance by a finite number
of arbitrarily near-sighted policemen.  (Some authors call this property \emph{precompact}.)
\end{defn}

\begin{prop} A metric space is compact if and only if its is complete and totally bounded.
\end{prop}

\begin{prop} A metric space is compact if and only if it is complete and totally bounded.
\end{prop}

\begin{defn} A subset $A$ of a topological space $X$ is
 \index{compact!relatively}%
 \index{relatively compact}%
\df{relatively compact} if its closure is compact.
\end{defn}

\begin{prop} In a metric space every relatively compact set is totally bounded.
\end{prop}

In a complete metric space the converse is true.

\begin{prop} In a complete metric space every totally bounded set is relatively compact.
\end{prop}

\begin{prop} In a metric space every totally bounded set is separable.
\end{prop}

\begin{thm}[Arzel\`a-Ascoli Theorem]\label{AAThm} Let $X$ be a compact Hausdorff space. A subset $\fml F$
of $\fml C(X)$ is totally bounded whenever it is (uniformly) bounded and equicontinuous.
\end{thm}

In the following we will be contrasting the \emph{weak} and \emph{strong} topologies on normed linear spaces.
The term
 \index{weak!\emph{vs.} strong topologies on normed spaces}%
 \index{topology!weak!\emph{vs.} strong on normed space}%
 \index{strong!\emph{vs.} weak topologies on normed spaces}%
 \index{conventions!on normed spaces \emph{strong topology} refers to the norm topology}%
\emph{weak topology} will refer to the usual weak topology determined by the bounded linear functionals
on the space as defined in~\ref{C067434}, while \emph{strong topology} will refer to the norm topology on the space.
Keep in mind, however, that the modifier ``strong'' in the context of normed linear spaces is, strictly speaking, 
redundant.  It is used solely for emphasis.  For example, many writers would state the next proposition as: 
\emph{Every weakly convergent sequence in a normed linear space is bounded.}

\begin{prop} Every weakly convergent sequence in a normed linear space is strongly bounded.
\end{prop}

\begin{proof}[\emph{Hint for proof}] Use the \emph{principle of uniform boundedness}~\ref{thm_PUB}. \ns
\end{proof}

\begin{defn} We say that a linear map $T\colon V \sto W$ is
 \index{weak!continuity}%
 \index{continuous!weakly}%
 \index{linear!map!weak continuity of a}%
\df{weakly continuous} if it takes weakly convergent nets to weakly convergent nets.
\end{defn}

\begin{prop} Every (strongly) continuous linear map between normed linear spaces is weakly continuous.
\end{prop}

\begin{defn} A linear map $T\colon V \sto W$ between normed linear spaces is
 \index{compact!linear map}%
 \index{linear!map!compact}%
 \index{operator!compact}%
\df{compact} if it takes bounded sets to relatively compact sets. Equivalently, $T$ is compact if it
maps the closed unit ball in $V$ onto a relatively compact set in~$W$.   The family of all compact linear maps
from $V$ to $W$ is denoted
 \index{K@$\ofml K(V,W)$!family of compact linear maps}%
 \index{K@$\ofml K(V)$!family of compact operators}%
by~$\ofml K(V,W)$; and we write $\ofml K(V)$ for~$\ofml K(V,V)$.
\end{defn}

\begin{prop} Every compact linear map is bounded.
\end{prop}

\begin{prop} A linear map $T\colon V \sto W$ between normed linear spaces is compact if and only if it maps
every bounded sequence in $V$ to a sequence in $W$ which has a convergent subsequence.
\end{prop}

\begin{exam} Every finite rank operator on a Banach space is compact.
\end{exam}

\begin{exam} Every bounded linear functional on a Banach space $B$ is a compact linear map from $B$ to~$\K$.
\end{exam}

\begin{exam} Let $C$ be the Banach space $\fml C([0,1])$ and $k \in \fml C([0,1] \times [0,1])$.  For each $f \in C$
 \index{integral!operator}%
 \index{operator!integral}%
define a function $Kf$ on $[0,1]$ by
   \[ Kf(x) = \int_0^1 k(x,y)f(y)\,dy\,. \]
Then the integral operator $K$ is a compact operator on~$C$.
\end{exam}

\begin{proof}[\emph{Hint for proof}] Use the \emph{Arzel\`a-Ascoli theorem}~\ref{AAThm}.  \ns
\end{proof}

\begin{exam}\label{X_square_int_kernel} Let $H$ be the Hilbert space $\lfs2{[0,1]}$ of (equivalence classes of) functions which
are square-integrable on $[0,1]$ with respect to Lebesgue measure and
 \index{integral!operator}%
 \index{operator!integral}%
$k$ be square-integrable function
on $[0,1] \times [0,1]$.  For each $f \in H$ define a function $Kf$ on $[0,1]$ by
   \[ Kf(x) = \int_0^1 k(x,y)f(y)\,dy\,. \]
Then the integral operator $K$ is a compact operator on~$H$.
\end{exam}

\begin{prop} An operator on a Hilbert space is compact if and only if it maps the closed unit ball in the space
onto a compact set.
\end{prop}

\begin{cor}\label{cor_id_op_not_cpt} The identity operator on an infinite dimensional Hilbert space is not compact.
(See example~\ref{X_l2ball_not_cpt}
\end{cor}

\begin{exam} If $B$ is a Banach space the family $\ofml B(B)$ of operators on $B$ is a
 \index{Banach!algebra!$\ofml B(B)$ as a}%
 \index{B@$\ofml B(V)$!as a unital Banach algebra}%
unital Banach algebra and the
 \index{K@$\ofml K(V)$!as a proper closed ideal}%
family $\ofml K(B)$ of compact operators on $B$ is a closed ideal in~$\ofml B(B)$.
If $B$ is infinite dimensional the ideal $\ofml K(H)$ is proper.
\end{exam}

\begin{prop} Let $T\colon B \sto C$ be a bounded linear map between Banach spaces. Then $T$ is compact
if and only if $T^*$ is.
\end{prop}

\begin{defn}\label{0013} A
 \index{C*@$C^*$-algebra}%
 \index{algebra!$C^*$-}%
\df{$C^\ast$-algebra} is a Banach algebra $A$ with involution which satisfies
   \[ \norm{a^*a} = \norm{a}^2 \]
for every $a \in A$.  This property of the norm is usually referred to as
 \index{C*@$C^*$-condition}%
\emph{the $C^*$-condition}. An algebra norm satisfying this condition is a
 \index{C*@$C^*$-norm}%
 \index{norm!$C^*$-}%
\df{$C^*$-norm}.  A
 \index{C*@$C^*$-subalgebra}%
 \index{subalgebra!$C^*$-}%
\df{$C^*$-subalgebra} of a $C^*$-algebra $A$ is a closed $*\,$-subalgebra of~$A$.

We denote by $\cat{CSA}$ the category of $C^*$-algebras and $*\,$-homomorphisms.
 \index{category!$\cat{CSA}$ as a}%
 \index{CSA@$\cat{CSA}$!the category}%
It may seem odd at first that in this category, whose objects have both algebraic and topological properties,
that the morphisms should be required to preserve only algebraic structure.  Remarkably enough, it turns out
that every morphism in $\cat{CSA}$ is automatically continuous---in fact, contractive (see
proposition~\ref{00152171})---and that
 \index{bijective morphisms!are invertible!in $\cat{CSA}$}%
every injective morphism is an isometry (see proposition~\ref{00152181}).  It is clear that
(as in definition~\ref{00109}) every bijective morphism in this category is an isomorphism.  And thus every bijective $*\,$-homomorphism
is an isometric $*\,$-isomorphism.
\end{defn}

\begin{exam}  The vector space $\C$ of complex numbers with the usual multiplication
 \index{C*@$C^*$-algebra!$\C$ as a}%
 \index{C@$\C$!as a $C^*$-algebra}%
of complex numbers and complex conjugation $z \mapsto \conj z$ as involution is a unital
commutative $C^*$-algebra.
\end{exam}

\begin{exam} If $X$ is a compact Hausdorff space, the algebra $\fml C(X)$ of continuous
 \index{C*@$C^*$-algebra!$\fml C(X)$ as a}%
 \index{C@$\fml C(X)$!as a $C^*$-algebra}%
complex valued functions on $X$ is a unital commutative $C^*$-algebra when involution is
taken to be complex conjugation.
\end{exam}

\begin{exam} If $X$ is a locally compact Hausdorff space, the Banach algebra
 \index{C*@$C^*$-algebra!$\fml C_0(X)$ as a}%
 \index{C@$\fml C_0(X)$!as a $C^*$-algebra}%
$\fml C_0(X) = \fml C_0(X,\C)$ of continuous complex valued functions on $X$ which vanish
at infinity is a (not necessarily unital) commutative $C^*$-algebra when involution is
taken to be complex conjugation.
\end{exam}

\begin{exam} If $(X,\mu)$ is a measure space, the algebra $\lfs\infty(X,\mu)$ of
 \index{C*@$C^*$-algebra!$L_\infty(S)$ as a}%
 \index{L@$\lfs\infty(S)$!as a $C^*$-algebra}%
essentially bounded measurable complex valued functions on $X$ (again with complex
conjugation as involution) is a $C^*$-algebra.  (Technically, of course, the members of
$\lfs\infty(X,\mu)$ are equivalence classes of functions which differ on sets of measure zero.)
\end{exam}

\begin{exam}\label{001305fa} The algebra $\ofml B(H)$ of bounded linear operators on a Hilbert space
 \index{C*@$C^*$-algebra!$\ofml B(H)$ as a}%
 \index{B@$\ofml B(V)$!as a $C^*$-algebra}%
$H$ is a unital $C^*$-algebra.  The family $\ofml K(H)$ of compact operators is a closed $*\,$-ideal
in the algebra.
\end{exam}

\begin{exam}\label{001306} As a special case of the preceding example note that the set
 \index{C*@$C^*$-algebra!M@$\mathbf M_n$ as a}%
 \index{Ma@$\mathbf M_n = \M n\C$!as a $C^*$-algebra}%
$\mathbf M_n$ of $n \times n$ matrices of complex numbers is a unital $C^*$-algebra.
\end{exam}

\begin{prop}\label{prop_K_is_clo_FR} In a Hilbert space $H$ the ideal $\ofml K(H)$ of compact operators is the closure of
the ideal $\ofml{FR}(H)$ of finite rank operators.
\end{prop}

In the decade or so before the Second World War, some leading Polish mathematicians gathered regularly in coffee
houses in Lwow to discuss results and propose problems.  First they met in the \emph{Caf\'e Roma} and later in the
\emph{Scottish Caf\'e}.  In 1935 Stefan Banach produced a bound notebook in which they could write down the problems.
This was the origin of the now famous \emph{Scottish Book}~\cite{Mauldin:1981}.  On November 6, 1936 Stanislaw Mazur
wrote down Problem 153, known as the \emph{approximation problem}, which asked, essentially, if the preceding proposition
holds for Banach spaces.  Can every compact Banach space operator be approximated in norm by finite rank operators?
Mazur offered as the prize for solving this problem a live goose (one of the more generous prizes promised in the
\emph{Scottish Book}).  The problem remained open for decades,  until 1972 when it was solved negatively by Per Enflo,
who succeeded in finding a separable reflexive Banach space in which such an approximation fails.  In a ceremony held
that year in Warsaw, Enflo received his live goose from Mazur himself.  The result was published in Acta
Mathematica~\cite{Enflo:1973} in~1973.

\begin{exam} The diagonal operator $\diag(1,\frac12,\frac13,\dots)$ is a compact operator on the Hilbert space~$l_2$.
\end{exam}











\section{Partial Isometries}
\begin{defn}  An element $v$ of a $C^*$-algebra $A$ is a
 \index{partial!isometry}%
\df{partial isometry} if $v^*v$ is a projection in~$A$.  (Since $v^*v$ is always
self-adjoint, it is enough to require that $v^*v$ be idempotent.)  The element $v^*v$ is
the
 \index{initial!projection}%
 \index{projection!initial}%
 \index{partial!isometry!initial projection of a}%
\df{initial} (or
 \index{support!projection}%
 \index{projection!support}%
 \index{partial!isometry!support projection of a}%
\df{support}) \df{projection} of~$v$ and $vv^*$ is the
 \index{final!projection}%
 \index{projection!final}%
 \index{partial!isometry!final projection of a}%
\df{final} (or
 \index{range!projection}%
 \index{projection!range}%
 \index{partial!isometry!range projection of a}%
\df{range}) \df{projection} of~$v$.  (It is an obvious consequence of the next
proposition that if $v$ is a partial isometry, then $vv^*$ is in fact a projection.)
\end{defn}

\begin{prop}\label{pi0111} If $v$ is a partial isometry in a $C^*$-algebra, then
 \begin{enumerate}[label=\rm{(\alph*)}]
  \item $vv^*v = v$\,; and
  \item $v^*$ is a partial isometry.
 \end{enumerate}
\end{prop}

\begin{proof}[\emph{Hint for proof}] Let $z = v - vv^*v$ and consider $z^*z$.  \ns
\end{proof}

\begin{prop}\label{pi0114} Let $v$ be a partial isometry in a $C^*$-algebra~$A$.  Then its initial
projection $p$ is the smallest projection (with respect to the partial ordering~$\preceq$
on $\fml P(A)$\,) such that $vp = v$ and its final projection $q$ is the smallest
projection such that $qv = v$.
\end{prop}

\begin{prop}\label{pi0117} If $V$ is a partial isometry on a Hilbert space $H$ (that is, if $V$ is a partial
isometry in the $C^*$-algebra~$\ofml B(H)$\,), then the initial projection $V^*V$ is the
projection of $H$ onto $(\ker V)^\perp$ and the final projection $VV^*$ is the projection
of $H$ onto~$\ran V$.
\end{prop}

Because of the preceding result $(\ker V)^\perp$ is called the
 \index{initial!space}%
 \index{space!initial}%
 \index{space!support}%
 \index{support!space}%
\df{initial} (or \df{support}) \df{space} of~$V$ and $\ran V$ is sometimes called the
 \index{final!space}%
 \index{space!final}%
\df{final space} of~$V$.)

\begin{prop}\label{pi0121} An operator on a Hilbert space is a partial isometry if and only if it is an
isometry on the orthogonal complement of its kernel.
\end{prop}





\begin{thm}[Polar Decomposition]\label{0022287} If $T$ is an operator on a Hilbert
 \index{polar decomposition}%
 \index{decomposition!polar}%
space $H$, then there exists a partial isometry $V$ on $H$ such that
 \begin{enumerate}
    \item[(i)] the initial space of $V$ is $(\ker T)^\perp$,
    \item[(ii)] the final space of $V$ is $\clo{\ran T}$, and
    \item[(iii)] $T = V\abs T$.
 \end{enumerate}
This decomposition of $T$ is unique in the sense that if $T = V_0P$ where $V_0$ is a
partial isometry and $P$ is a positive operator such that $\ker V_0 = \ker P$, then $P =
\abs T$ and $V_0 = V$.
\end{thm}

\begin{proof} See \cite{Conway:1990}, VIII.3.11; \cite{Davidson:1996}, theorem I.8.1;
\cite{KadisonR:1983}, theorem 6.1.2; \cite{Murphy:1990}, theorem 2.3.4;
or~\cite{Pedersen:1995}, theorem 3.2.17. \ns
\end{proof}

\begin{cor} If $T$ is an
  \index{polar decomposition!of an invertible operator}%
 \index{operator!invertible!polar decomposition of}%
 \index{invertible!operator!polar decomposition of}%
invertible operator on a Hilbert space, then the partial isometry in the polar
decomposition (see~\ref{0022287}) is unitary.
\end{cor}

\begin{prop} If $T$ is a
 \index{polar decomposition!of a normal operator}%
 \index{operator!normal!polar decomposition of}%
 \index{normal!operator!polar decomposition of}%
normal operator on a Hilbert space $H$, then there exists a unitary operator $U$ on $H$
which commutes with $T$ and satisfies $T = U\abs T$.
\end{prop}

\begin{proof} See \cite{Pedersen:1995}, proposition 3.2.20.\ns \end{proof}

\begin{exer} Let $(S, \fml A,\mu)$ be a $\sigma$-finite measure space and $\phi \in
L_\infty(\mu)$. Find the polar decomposition of the
 \index{multiplication operator!polar decomposition of}%
 \index{operator!multiplication!polar decomposition of}%
 \index{polar decomposition!of a multiplication operator}%
multiplication operator $M_\phi$ (see example~\ref{C063527}).
\end{exer}









\section{Trace Class Operators}

Let's start by reviewing some of the standard properties of the \emph{trace} of an $n \times n$ matrix.

\begin{defn}\label{tr0113} Let $a = [a_{ij}] \in \mathbf M_n$. Define $\tr a$, the
 \index{trace!of a matrix}%
 \index{matrix!trace of a}%
 \index{tr@$\tr$ (trace)}%
\df{trace} of $a$ by
  \[ \tr a = \sum_{i = 1}^n a_{ii}. \]
\end{defn}

\begin{prop}\label{tr0117} The function $\tr\colon \mathbf M_n \sto \C$ is linear.
\end{prop}

\begin{prop}\label{tr0121} If $a$, $b \in \mathbf M_n$, then $\tr (ab) = \tr (ba)$.
\end{prop}

The trace is invariant under similarity.  Two $n \times n$ matrices $a$ and $b$ are
 \index{similar!matrices}%
 \index{matrices!similar}%
\df{similar} if there exists an invertible matrix $s$ such that $b = sas^{-1}$.

\begin{prop}\label{tr0124} If matrices in $\mathbf M_n$ are similar, then they have the same trace.
\end{prop}

\begin{defn}\label{tr0211} Let $H$ be a separable Hilbert space, $(e_n)$ be an orthonormal basis for~$H$, and $T$ be
a positive operator on~$H$.  Define the
 \index{trace}%
 \index{tr@$\tr$ (trace)}%
\df{trace} of~$T$~by
    \[ \tr T:=\sum_{k=1}^\infty \langle Te_k,e_k \rangle. \]
(The trace is not necessarily finite.)
\end{defn}

\begin{prop} If $S$ and $T$ are positive operators on a separable Hilbert space and $\alpha \in  \K$, then
   \[  \tr(S + T) = \tr S + \tr T \]
and
   \[ \tr(\alpha T) = \alpha \tr T. \]
\end{prop}

\begin{prop} If $T$ is an operator on a separable Hilbert space then $\tr(T^*T) = \tr(TT^*)$.
\end{prop}

To show that the definition given above of the trace of a positive operator does not depend on the orthonormal basis chosen for the
Hilbert space we will need a result (given next) which depends on a fact which appears later in these notes: every positive Hilbert
space operator has a (positive) square root (see example~\ref{X_sqroot_op}).

\begin{prop}\label{tr0221} On the family of positive operators on a separable Hilbert space $H$ the trace is invariant under unitary equivalence;
that is, if $T$ is positive and $U$ is unitary, then $\tr T = \tr(UTU^*)$.
\end{prop}

\begin{proof}[\emph{Hint for proof}] In the preceding proposition let $S = UT^{\frac12}$.    \ns
\end{proof}

\begin{prop}\label{tr0224} The definition of the \emph{trace}~\ref{tr0211} of a positive operator on a separable Hilbert space is independent
of the choice of orthonormal basis for the space.
\end{prop}

\begin{proof}[\emph{Hint for proof}] Let $(e^k)$ and $(f^k)$ be orthonormal bases for the space and $T$ be a positive operator
on the space.  Take $U$ to be the unique unitary operator with the property that $e^k = Tf^k$ for each $k \in \N$.  \ns
\end{proof}

\begin{defn} Let $V$ be a vector space.  A subset $C$ of $V$ is a
 \index{cone}%
\df{cone} in $V$ is $\alpha C \subseteq C$ for every $\alpha \ge 0$.  A cone $C$ in $V$ is
 \index{cone!proper}%
 \index{proper!cone}%
\df{proper} is $C \cap (-C) = \{\vc 0\}$.
\end{defn}

 \begin{prop} A cone $C$ in a vector space is convex if and only if $C + C \subseteq C$.
\end{prop}

\begin{exam} On a separable Hilbert space the family of all positive operators with finite trace forms a proper convex cone in the vector
space~$\ofml B(H)$.
\end{exam}

\begin{defn} In a separable Hilbert space $H$ the linear subspace of $\ofml B(H)$ spanned by the positive operators with finite trace
is the family $\ofml T(H)$ of
 \index{trace!class}%
 \index{operator!trace class}%
 \index{T@$\ofml T(H)$ (trace class operators on~$H$)}%
\df{trace class} operators on~$H$.
\end{defn}

\begin{prop} Let $H$ be a separable Hilbert space with orthonormal basis~$(e_n)$.  If $T$ is a trace class operator on $H$, then
   \[ \tr T = \sum_{k=1}^\infty \langle Te_k,e_k \rangle \]
and the sum on the right is absolutely convergent.
\end{prop}
















\section{Hilbert-Schmidt Operators}

\begin{defn} An operator $A$ on a separable Hilbert space $H$ is a
 \index{Hilbert!-Schmidt operator}%
 \index{operator!Hilbert-Schmidt}%
\df{Hilbert-Schmidt} operator if $A^*A$ is of trace class; that is, if $\sum_{k=1}^\infty \norm{Ae_k}^2 < \infty$,
where $(e_k)$ is an orthonormal basis for~$H$.  The family of all Hilbert-Schmidt operators on $H$ is
 \index{HS@$\ofml{HS}(H)$ (Hilbert-Schmidt operators on~$H$)}%
denoted by~$\ofml{HS}(H)$.
\end{defn}

\begin{prop} If $H$ is a separable Hilbert space, then the family of Hilbert-Schmidt operators on $H$ is a two-sided
ideal in~$\ofml B(H)$.
\end{prop}

\begin{exam}  The family of Hilbert-Schmidt operators on a Hilbert space $H$ can be made into an inner product space.
\end{exam}

\begin{proof}[\emph{Hint for proof}]  Polarization.  If $A$, $B \in \ofml{HS}(H)$,
consider $\sum_{k=0}^3 i^k \bigl(A + i^kB\bigr)^*(A + i^kB)$.  \ns
\end{proof}

\begin{exam} The inner product you defined in the preceding example makes $\ofml{HS}(H)$ into a Hilbert space.
\end{exam}

\begin{prop} On a separable Hilbert space every Hilbert-Schmidt operator is compact.
\end{prop}

\begin{proof}[\emph{Hint for proof}] If $(e_k)$ is an orthonormal basis for the Hilbert space, consider the
operators $F_n := AP_n$ where $P_n$ is the projection onto the span of $\{e_1 \dots, e_n\}$. \ns
\end{proof}

\begin{exam} The integral operator defined in example~\ref{X_square_int_kernel} is a Hilbert-Schmidt operator.
\end{exam}

\begin{exam}
    The Volterra operator defined in example~\ref{000319} is a Hilbert-Schmidt operator.
\end{exam}











\endinput
